\documentclass[12pt,a4paper]{article}

\usepackage[T2A]{fontenc}
\usepackage[utf8]{inputenc}
\usepackage[russian]{babel}
\usepackage{amsmath, amsfonts, amssymb, amsthm}
\usepackage{geometry}
\geometry{top=2cm, bottom=2cm, left=2.5cm, right=1.5cm}
\usepackage{verbatim}
\usepackage{hyperref}
\usepackage{xcolor}

\hypersetup{
    colorlinks=true,
    linkcolor=blue,
    urlcolor=blue
}

\newtheorem{definition}{Определение}
\newtheorem{theorem}{Теорема}
\newtheorem{lemma}{Лемма}
\newtheorem{proof_custom}{Доказательство}

\title{Билет №8 \\ \small Структуры данных. Линейные (списки, очереди, деки, вектора). Очереди с приоритетами. Деревья поиска. (ИТМО)}
\author{Сериков Владислав Александрович}
\date{21 апреля 2026}

\begin{document}

\maketitle
\tableofcontents
\newpage

\section{Введение: Асимптотический анализ}

\begin{definition}[$\mathcal{O}$-большое]
$f(n) = \mathcal{O}(g(n))$, если существуют такие константы $c > 0$ и $n_0 > 0$, что для всех $n \ge n_0$ выполняется $0 \le f(n) \le c \cdot g(n)$. Обозначает верхнюю границу (наихудший случай).
\end{definition}

\begin{definition}[$\Omega$-большое и $\Theta$-большое]
$f(n) = \Omega(g(n))$ задает нижнюю границу алгоритма. Если $f(n) = \mathcal{O}(g(n))$ и $f(n) = \Omega(g(n))$, то $f(n) = \Theta(g(n))$, что означает точную асимптотическую оценку.
\end{definition}

\section{Линейные структуры данных}

\subsection{Векторы (Динамические массивы)}
\begin{definition}
\textbf{Вектор} — это структура данных, представляющая собой массив, способный автоматически изменять свой размер (capacity) при добавлении новых элементов.
\end{definition}

\textbf{Алгоритм вставки в конец (push\_back):}
\begin{enumerate}
    \item Если текущее количество элементов $size$ равно вместимости $capacity$:
    \begin{enumerate}
        \item Выделить новый участок памяти размера $capacity \times 2$.
        \item Скопировать $size$ элементов из старого массива в новый.
        \item Освободить старую память.
    \end{enumerate}
    \item Записать новый элемент в индекс $size$.
    \item $size \leftarrow size + 1$.
\end{enumerate}

\textbf{Минимальный пример:}
Изначально: \texttt{arr = [5, 6]}, $size = 2$, $capacity = 2$.
Вызов \texttt{push\_back(7)}:
1. Массив переполнен. Аллоцируем новый массив размера 4.
2. Копируем старые: \texttt{[5, 6, \_, \_]}.
3. Вставляем 7: \texttt{[5, 6, 7, \_]}. $size = 3$.

\begin{theorem}
Амортизированная сложность операции \texttt{push\_back} составляет $\mathcal{O}(1)$.
\end{theorem}
\begin{proof}
Применим метод группировки (aggregate analysis). Пусть мы делаем $N$ операций \texttt{push\_back}. Обозначим $c_i$ — реальную стоимость $i$-й операции.
Если $i-1$ является степенью двойки, происходит реаллокация и копирование $i-1$ элементов, тогда $c_i = i$. В противном случае $c_i = 1$.
Суммарное время $T(N)$ для $N$ вставок:
$$ T(N) = \sum_{i=1}^{N} c_i \le N + \sum_{j=0}^{\lfloor \log_2 N \rfloor} 2^j $$
Используя формулу суммы геометрической прогрессии:
$$ \sum_{j=0}^{\lfloor \log_2 N \rfloor} 2^j = 2^{\lfloor \log_2 N \rfloor + 1} - 1 \le 2N $$
Тогда $T(N) \le 3N$. Средняя (амортизированная) стоимость одной операции: $\frac{T(N)}{N} \le 3 = \mathcal{O}(1)$.
\end{proof}

\subsection{Связные списки (Linked Lists)}
\begin{definition}
\textbf{Связный список} — линейная структура, состоящая из узлов. Каждый узел содержит полезные данные и указатель на следующий (а в двусвязном — и на предыдущий) узел.
\end{definition}

\textbf{Алгоритм вставки узла B после узла A (односвязный список):}
\begin{enumerate}
    \item Создать узел B с нужным значением.
    \item \texttt{B.next = A.next}
    \item \texttt{A.next = B}
\end{enumerate}

\textbf{Иллюстрация:}
\begin{verbatim}
Изначально: (A) -> (C)
Шаг 1, 2:   (B) -> (C)
Шаг 3:      (A) -> (B) -> (C)
\end{verbatim}
Вставка при известном указателе работает за $\mathcal{O}(1)$, но поиск узла требует $\mathcal{O}(n)$.

\subsection{Очереди и Деки}
\begin{definition}
\textbf{Очередь (Queue)} — коллекция элементов, реализующая принцип FIFO (First-In, First-Out).
\textbf{Дек (Deque, Double-Ended Queue)} — двусторонняя очередь, поддерживающая вставку и удаление элементов как с начала, так и с конца за $\mathcal{O}(1)$.
\end{definition}

Реализуются либо поверх двусвязных списков, либо на кольцевых буферах (циклических массивах).
\textbf{Алгоритм Enqueue (вставка в кольцевой буфер):}
\begin{enumerate}
    \item \texttt{buffer[tail] = x}
    \item \texttt{tail = (tail + 1) mod capacity}
\end{enumerate}

\section{Очереди с приоритетами}
\begin{definition}
\textbf{Очередь с приоритетами} — абстрактный тип данных, поддерживающий операции \texttt{insert} и \texttt{extract\_min} (или \texttt{extract\_max}).
Чаще всего реализуется структурой \textbf{Бинарная куча (Binary Heap)}.
\end{definition}

\begin{definition}[Min-Heap]
Бинарное дерево, удовлетворяющее двум свойствам:
\begin{enumerate}
    \item \textbf{Структурное:} Почти полное бинарное дерево (все уровни заполнены, кроме, возможно, последнего, который заполнен слева направо).
    \item \textbf{Свойство кучи:} Для любого узла $v$, отличного от корня, $key(parent(v)) \le key(v)$.
\end{enumerate}
\end{definition}

\textbf{Алгоритм вставки (\texttt{insert(x)}):}
\begin{enumerate}
    \item Добавить $x$ в конец массива, представляющего кучу.
    \item Вызвать функцию \texttt{sift\_up} (всплытие) для нового элемента.
\end{enumerate}

\textbf{Шаги \texttt{sift\_up(i)}:}
Пока $i > 0$ и элемент $A[i] < A[parent(i)]$ (где $parent(i) = \lfloor (i-1)/2 \rfloor$):
меняем их местами и обновляем $i \leftarrow parent(i)$.

\textbf{Пример и иллюстрация \texttt{insert(1)} в кучу [2, 3, 4]:}
\begin{verbatim}
1. Начальная куча      2. Добавляем 1        3. sift_up(3)        4. sift_up(1)
(массив [2, 3, 4]):    (массив [2, 3, 4, 1]) Сравниваем 1 и 3     Сравниваем 1 и 2
      2                      2                      2                    1
     / \                    / \                    / \                  / \
    3   4                  3   4                  1   4                2   4
                          /                      /                    /
                         1                      3                    3
\end{verbatim}

\begin{theorem}[Корректность sift\_up]
После выполнения \texttt{sift\_up} массив снова является корректной Min-Heap.
\end{theorem}
\begin{proof}
Добавив элемент в конец, мы сохранили структурное свойство дерева. Единственное возможное нарушение свойства кучи находится между новым узлом $i$ и его родителем $p$.
Если $A[i] < A[p]$, мы меняем их местами. Теперь новый родитель $A[p]$ (бывший $A[i]$) строго меньше старого родителя, а значит, он также меньше второго своего потомка (брата узла $i$, если тот есть). Свойство кучи в данном поддереве восстановлено. Нарушение могло сместиться только на уровень выше. Поскольку высота дерева равна $\lfloor \log_2 n \rfloor$, процесс завершится максимум за $\mathcal{O}(\log n)$ шагов, достигнув корня, либо когда $A[i] \ge A[p]$. Инвариант кучи восстанавливается.
\end{proof}

\section{Деревья поиска (Binary Search Trees)}

\begin{definition}
\textbf{Бинарное дерево поиска (BST)} — бинарное дерево, обладающее следующим инвариантом: для любого узла $u$ все ключи в его левом поддереве строго меньше $u.key$, а все ключи в правом поддереве строго больше $u.key$.
\end{definition}

\textbf{Алгоритм поиска (\texttt{search(x, node)}):}
\begin{enumerate}
    \item Если \texttt{node} равен \texttt{null} или \texttt{node.key == x}, вернуть \texttt{node}.
    \item Если \texttt{x < node.key}, вернуть \texttt{search(x, node.left)}.
    \item Иначе вернуть \texttt{search(x, node.right)}.
\end{enumerate}

\begin{theorem}
Алгоритм поиска в BST корректен и находит элемент $x$ тогда и только тогда, когда $x$ присутствует в дереве.
\end{theorem}
\begin{proof}
Доказательство по индукции. База: если узел пуст, элемента в дереве нет.
Переход: пусть мы находимся в узле $u$. Согласно инварианту BST, элементы левого поддерева $< u.key$, а правого $> u.key$. Следовательно, если искомый $x < u.key$, он ни при каких условиях не может находиться в правом поддереве или быть равным $u.key$. Мы безопасно отсекаем правое поддерево и продолжаем поиск в левом. Транзитивность отношения порядка гарантирует, что мы не пропустим $x$.
\end{proof}

\textbf{Алгоритм вставки (\texttt{insert(x)}):}
\begin{enumerate}
    \item Спускаемся по дереву аналогично алгоритму \texttt{search}.
    \item Как только должны перейти в \texttt{null}-потомка, создаём на этом месте новый узел с ключом $x$.
\end{enumerate}

\textbf{Минимальный пример вставки ключа 4:}
\begin{verbatim}
Изначальное BST:       Шаг 1: 4 > 3      Шаг 2: 4 < 5        Результат:
      3                 (идем вправо)    (идем влево)             3
     / \                                                         / \
    1   5                                                       1   5
                                                                   /
                                                                  4
\end{verbatim}

\textbf{Асимптотика BST (с отсылкой на статью-шпаргалку Хабра):}
В среднем случае операции поиска, вставки и удаления работают за $\mathcal{O}(\log n)$, так как высота дерева пропорциональна логарифму числа элементов.
Однако в худшем случае (вставка уже отсортированной последовательности) дерево вырождается в связный список (высота $\mathcal{O}(n)$). Чтобы этого избежать, применяются \textbf{сбалансированные} деревья поиска (АВЛ-дерево, Красно-черное дерево), где высота поддерживается строгими инвариантами и вращениями.

\subsection{Самобалансирующиеся деревья поиска}

Как было отмечено ранее, наивное бинарное дерево поиска может вырождаться, давая время выполнения $\mathcal{O}(n)$. Чтобы гарантировать логарифмическую сложность в худшем случае (или в амортизированном смысле), вводятся дополнительные структурные инварианты.

\subsubsection{АВЛ-дерево (AVL Tree)}
\begin{definition}
\textbf{АВЛ-дерево} — бинарное дерево поиска, в котором для любого узла высота его левого и правого поддеревьев отличается не более чем на 1. Разность высот называется баланс-фактором (balance factor): $bf(node) \in \{-1, 0, 1\}$.
\end{definition}

Основной механизм восстановления баланса — \textbf{вращения (rotations)}. Вращение меняет структуру дерева локально за $\mathcal{O}(1)$, сохраняя инвариант BST (порядок ключей).

\textbf{Иллюстрация правого и левого вращения:}
\begin{verbatim}
      y                                 x
     / \        Right Rotate(y)        / \
    x   T3     ---------------->      T1  y
   / \         <----------------         / \
  T1  T2        Left Rotate(x)          T2  T3
\end{verbatim}

\textbf{Алгоритм балансировки после вставки узла:}
\begin{enumerate}
    \item Обновить высоты предков вставленного узла.
    \item Если для какого-то узла $|bf| > 1$, применить вращения.
    \begin{enumerate}
        \item Если $bf > 1$ и вставка была в левое поддерево левого ребенка: Малое правое вращение.
        \item Если $bf > 1$ и вставка была в правое поддерево левого ребенка: Большое правое вращение (сначала левое вращение ребенка, затем правое вращение родителя).
        \item Симметричные случаи для $bf < -1$.
    \end{enumerate}
\end{enumerate}

\begin{theorem}[Оценка высоты АВЛ-дерева]
Высота АВЛ-дерева из $n$ узлов ограничена сверху $\mathcal{O}(\log n)$.
\end{theorem}
\begin{proof}
Рассмотрим минимальное количество узлов $N(h)$ в АВЛ-дереве высоты $h$. Чтобы минимизировать число узлов при заданной высоте, дерево должно быть максимально несбалансированным. Одно поддерево корня должно иметь высоту $h-1$, а другое $h-2$. Тогда:
$$ N(h) = 1 + N(h-1) + N(h-2) $$
Это рекуррентное соотношение подобно числам Фибоначчи. Заметим, что $N(h) > 2 \cdot N(h-2)$. Применяя это неравенство по индукции, получаем $N(h) \ge 2^{h/2}$.
Отсюда $h \le 2 \log_2 N(h)$. Таким образом, $h = \mathcal{O}(\log n)$, что строго доказывает логарифмическое время поиска.
\end{proof}

\subsubsection{Красно-чёрное дерево (Red-Black Tree)}
\begin{definition}
\textbf{Красно-чёрное дерево} — бинарное дерево поиска, каждому узлу которого приписан цвет (красный или чёрный), удовлетворяющее пяти свойствам:
\begin{enumerate}
    \item Каждый узел либо красный, либо чёрный.
    \item Корень дерева — чёрный.
    \item Все листья (NIL-узлы, не содержащие данных) — чёрные.
    \item У красного узла оба потомка — чёрные (недопустимы два красных узла подряд).
    \item Для любого узла все простые пути от него до листьев, являющихся его потомками, содержат одинаковое количество чёрных узлов (это число называется чёрной высотой, $bh(x)$).
\end{enumerate}
\end{definition}

\textbf{Алгоритм вставки (кратко):}
Новый узел всегда вставляется как красный. Если его родитель также красный (нарушение свойства 4), происходит перекраска узлов (recoloring) или вращение. Действия зависят от цвета "дяди" (брата родителя).

\begin{theorem}
Красно-чёрное дерево с $n$ внутренними узлами имеет высоту $h \le 2 \log_2(n+1)$.
\end{theorem}
\begin{proof}
Сначала докажем лемму: поддерево с корнем $x$ содержит не менее $2^{bh(x)} - 1$ внутренних узлов. (Доказывается по индукции по высоте узла).
Теперь рассмотрим корень дерева $root$. По свойству 4 (отсутствие двух подряд идущих красных узлов) на любом пути от корня до листа как минимум половина узлов — чёрные. Следовательно, чёрная высота корня $bh(root) \ge h/2$.
Применяя лемму к корню:
$$ n \ge 2^{bh(root)} - 1 \ge 2^{h/2} - 1 $$
$$ n + 1 \ge 2^{h/2} \implies \log_2(n+1) \ge \frac{h}{2} \implies h \le 2 \log_2(n+1) $$
Из теоремы следует, что операции \texttt{search}, \texttt{insert} и \texttt{delete} в худшем случае работают за $\mathcal{O}(\log n)$.
\end{proof}

\subsubsection{Сплей-дерево (Splay Tree)}
\begin{definition}
\textbf{Сплей-дерево} (расширяющееся дерево) — самобалансирующееся бинарное дерево поиска без строгих структурных инвариантов высоты, которое использует эвристику: недавно запрошенные элементы должны быть легко доступны. Основано на операции \texttt{splay}.
\end{definition}

\textbf{Алгоритм \texttt{Splay(x)}:}
Операция перемещает узел $x$ в корень дерева с помощью серии специфических вращений. Пока $x$ не корень, выполняются шаги:
\begin{itemize}
    \item \textbf{Zig:} Если родитель $p$ узла $x$ является корнем. Выполняется одно вращение.
    \item \textbf{Zig-Zig:} Если $x$ и $p$ оба являются левыми (или оба правыми) потомками своих родителей. Сначала вращается родитель $p$, затем сам $x$. (Отличие от обычного вращения к корню!).
    \item \textbf{Zig-Zag:} Если $x$ — правый потомок левого родителя (или наоборот). Вращается $x$, а затем снова $x$.
\end{itemize}

\begin{theorem}
Амортизированная стоимость любой операции (поиск, вставка, удаление) в Splay-дереве составляет $\mathcal{O}(\log n)$.
\end{theorem}
\begin{proof}[Идея доказательства (Метод потенциалов Слеатора и Тарьяна)]
Каждому узлу $x$ присваивается вес (обычно 1). Пусть $S(x)$ — сумма весов всех узлов в поддереве с корнем $x$. Определим ранг узла как $r(x) = \log_2 S(x)$. Потенциал всего дерева $\Phi(T) = \sum_{x \in T} r(x)$.
Амортизированное время операции равно $t_{amortized} = t_{real} + \Delta \Phi$.
Анализ каждого шага (zig, zig-zig, zig-zag) показывает, что амортизированная стоимость одного перемещения узла $x$ к корню ограничена сверху $3(r(root) - r(x)) + 1$. Поскольку $r(root) \le \log_2 n$, общая амортизированная сложность операции составляет $\mathcal{O}(\log n)$.
\end{proof}

\end{document}
